\documentclass[10pt,a4paper]{amsart}
\usepackage[margin=19mm]{geometry}
\usepackage{amsmath,amssymb,mathtools}
\usepackage[scr=rsfs,cal=euler]{mathalfa}
\usepackage{xfrac}
\usepackage[unicode,hidelinks,bookmarks=false]{hyperref}
\newcommand{\ie}{i.e.\ }
\newcommand{\cf}{cf.\ }
\newcommand{\resp}{respectively\ }
\newcommand{\Subsection}{\subsection}
\providecommand{\nobreakdash}{\nobreak\hbox{-}}
\setlength{\emergencystretch}{2em}
\multlinegap0pt
\usepackage{bm}
\usepackage{bbm}
\usepackage{dsfont}
\usepackage{xspace}
\usepackage{cancel}
\usepackage{mathtools}
\usepackage{amsthm}
\makeatletter
\@ifundefined{theorem}{\newtheorem{theorem}{Theorem}}{}
\@ifundefined{corollary}{\newtheorem{corollary}[theorem]{Corollary}}{}
\makeatother
\title[Regularity of the free boundary in an unstable parabolic pro]{Regularity of the free boundary in an unstable parabolic problem}
\author{Mark Allen, Gilles Bokolo-Tamba}
\date{Source: arXiv:2508.15997v1 (CC BY 4.0)}
\begin{document}
\maketitle

\noindent\emph{Source and licence.} Regularity of the free boundary in an unstable parabolic problem. Version arXiv:2508.15997v1 (CC BY 4.0), distributed under the Creative Commons Attribution 4.0 International licence, as recorded in the arXiv licence metadata for this exact version and shown on the abstract page https://arxiv.org/abs/2508.15997v1. Full rights notice: licenses/arxiv-2508.15997-CC-BY-4.0.txt.\par\smallskip

\noindent\textit{Editorial scope.} Counted unit: the Theorem whose source label is \texttt{t:holder} in the article (source0.tex lines 764--766, source label \texttt{t:holder}), delivered together with the article's own complete original proof of it (source0.tex lines 768--867, one \texttt{proof} environment of 100 lines). No mathematics is written, repaired or re-derived: statement and proof are the article's own text line for line. Required context retained uncounted: e:main (lines 46--48); c:rescale (lines 619--628). Every definition, display and label of those lines is retained unchanged. Omitted from this package and not redistributed: 1--45, 49--618, 629--763, 767--767, 868--1123. Nothing inside a retained excerpt is omitted or altered: every retained body is delivered exactly as the article's lines are written, so no omission receipt of any kind accompanies this package. Citations: the retained excerpts carry no citation command at all, so no bibliography entry of the article is transcribed and no reference is redistributed. Reference adaptation: none was needed and none was made; every reference inside the delivered unit resolves inside it, so no unresolved reference or citation is produced. Printed slips kept literal: no printed word, symbol, hypothesis or number of the counted statement or of its proof was altered. Layout and class: the article's own class is \texttt{amsart} with packages including graphicx, amsmath, amssymb, amsthm; the wrapper is the lane's pinned \texttt{amsart} editorial wrapper at 10pt on A4, which restates the theorem-like environments the retained text needs and adds the article's own macro definitions that the retained text needs (none), with extra wrapper packages (bm, bbm, dsfont, xspace, cancel, mathtools). The delivered numbering is the unit's own. Assets: no retained span contains a figure environment, a graphics-inclusion command, a PDF-page inclusion, a table or a TikZ picture; no image file is copied into this package, the package declares no asset dependency and no image file is redistributed with it. Licence: version arXiv:2508.15997v1 of this article is distributed under the Creative Commons Attribution 4.0 International licence, as recorded in the arXiv licence metadata for this exact version and shown on the abstract page https://arxiv.org/abs/2508.15997v1. A full rights notice is delivered at licenses/arxiv-2508.15997-CC-BY-4.0.txt. Characters: every retained excerpt is pure ASCII, so the delivered engine, XeLaTeX with its default Latin Modern text font, reports no missing character.
 \begin{equation} \label{e:main}
     u_t - \Delta u = \chi_{\{u>0\}}, 
 \end{equation}

\begin{corollary} \label{c:rescale}
 Assume $u$ is a solution to \eqref{e:main} in $Q_1$. Fix $0<\gamma<1-\alpha$, and assume that $|\nabla u(x,t)|>0$. Let $\rho$ be defined as $\rho^{1-\gamma}=|\nabla u(x,t)|/M$. There exists a constant $C_1$ depending on $\gamma$ such that 
 \begin{equation} \label{e:rho}
 \sup_{Q_{\rho/2}(x,t)} |\partial_t u| \leq C_1 M\rho^{-\gamma},
 \end{equation}
 and a constant $C_2$ depending on $\alpha,\gamma$ such that  
 \begin{equation} \label{e:rho2}
  | u_t(y,s)-u_t(x,t)| \leq C_2 r^{-2\gamma-\alpha} (|x-y|^{\alpha}+|t-s|^{\alpha/2}) \quad \text{ for any } (y,s) \in Q_{r/2}. 
 \end{equation}
\end{corollary}

\begin{theorem} \label{t:holder}
    The free boundary $\Gamma$ is a $C^{1,\alpha/2}$ function of the spatial variable $x$ for any $0<\alpha<1$. 
\end{theorem}

\begin{proof}
    Fix $0<\alpha<1$, and let $\alpha<\beta<1$. Fix $(x,t) \in \Gamma \cap Q_{1/2}$. If $|\nabla u(x,t)|=0$, then 
    $|\nabla u(y,s)|\leq |x-y|^{\alpha}+|t-s|^{\alpha/2}$, so that the angle $\theta$ between $\nu_{x,t}$ and $\nu_{y,s}$ satisfies
    \[
    \sin \theta \leq \frac{|\nabla u(y,s)|}{\sqrt{|\nabla u(y,s)|^2 + c^2}}
    \leq \frac{1}{c} (|x-y|^{\alpha/2}+|t-s|^{\alpha/2}). 
    \]
    We now assume that $|\nabla u(x,t)|>0$. 
    
   \noindent  \textbf{Case 1: $|x-y|+|t-s|^{1/2}\geq r/2$.}

     The angle $\theta$ between $\nu_{x,t}$ and $(0,1)$ is given by 
    \[
     \begin{aligned}
     \sin \theta_k&=\frac{|\nabla u(x,t)|}{\sqrt{|\nabla u(x,t)|^2 + (\partial_t u(x,t))^2}} \\
    &\leq \frac{|\nabla u(x,t)|}{c} \\
    &= \frac{Mr^{\beta}}{c}\\
    &\leq \frac{4M}{c}\left(|x - y|^{\beta}+|t-s|^{\beta/2} \right)\\
    &\leq \frac{4M}{c}\left(|x - y|^{\beta/2}+|t-s|^{\beta/2} \right) \\
    \end{aligned}
    \]
    Now
    \[
    \begin{aligned}
    |\nabla u(y,s)|&\leq |\nabla u(x,t)|+|\nabla u(x,t)-\nabla u(y,s)| \\
    &\leq Mr^{\beta} + |x-y|^{\beta}+ |t-s|^{\beta/2}\\
    &\leq 5M(|x-y|^{\beta/2}+ |t-s|^{\beta/2}).
    \end{aligned}
    \]
    Thus, a similar argument gives that the angle $\theta$ between $\nu_{y,s}$ and $(0,1)$ satisfies
    \[
    \sin \theta \leq \frac{10M}{c}\left(|x - y|^{\beta/2}+|t-s|^{\beta/2} \right).
    \]


   \noindent \textbf{Case 2: $|x-y|+|t-s|^{1/2}\leq r/2$.}
       To relate to the normal, we use the following vector valued function which projects onto the sphere. 
    \[
    F(\xi):= \frac{\xi}{|\xi|},
    \]
    for $\xi \in \mathbb{R}^{n+1}$. We will let $\xi_j =u_{x_j}(y,s)$ for $1\leq j \leq n$ and $\xi_{n+1}=u_t(y,s)$. If the differences $u_{x_i}(x,t)-u_{x_i}(y,s)$ and $u_t(x,t)-u_t(y,s)$ are small, then
    \[
     \nu_{x,t}-\nu_{y,s}\approx
     \begin{bmatrix}
         \frac{\partial F^1}{\partial \xi_1} & \dots &\frac{\partial F^1}{\partial \xi_n}
         & \frac{\partial F^1}{\partial \xi_{n+1}} \\
         \frac{\partial F^2}{\partial \xi_1} & \dots &\frac{\partial F^2}{\partial \xi_n}
         & \frac{\partial F^2}{\partial \xi_{n+1}} \\
          \vdots &\ddots &\vdots & \vdots\\
         \frac{\partial F^{n+1}}{\partial \xi_1} & \dots &\frac{\partial F^{n+1}}{\partial \xi_n}
         & \frac{\partial F^{n+1}}{\partial \xi_{n+1}} \\
     \end{bmatrix}
     \begin{bmatrix}
         u_{x_1}(x,t)-u_{x_1}(y,s) \\
         u_{x_2}(x,t)-u_{x_2}(y,s) \\
         \vdots \\
         \\
         u_{t}(x,t)-u_{t}(y,s) \\
     \end{bmatrix}.
    \]
    Now the difference $|u_{x_i}(x,t)-u_{x_i}(y,s)|$ will be small from H\"older continuity of the gradient, but $|u_t(x,t)-u_t(y,s)|$ is not necessarily small. However, we now show that we can still obtain a bound. First, 
    \[
    \frac{\partial F^j}{\partial \xi_k} = \frac{\delta_{kj}|\xi|^2 - \xi_k \xi_j}{|\xi|^{3/2}}. 
    \]
    Notice that $|\partial_{\xi_k} F^j|\leq 1/c$ since $|\eta| \geq c$. If $k \leq n$, then 
    $|u_{x_i}(x,t)-u_{x_i}(y,s)|$ is small, and 
    \[
     \left|\frac{\partial F^j}{\partial \xi_k} (u_{x_k}(x,t) -u_{x_k}(y,s))\right|
     \leq c^{-1}|(u_{x_k}(x,t) -u_{x_k}(y,s))| \leq c^{-1}(|x-y|^{\alpha} + |t-s|^{\alpha/2}). 
    \]
    If $k=j=n+1$, then $F^{n+1}$ is concave as a function of $\xi_{n+1}$. Then 
    \[
    |F^{n+1}(\xi_1,\ldots,\xi_n,u_t(x,t))- F^{n+1}(\xi_1,\ldots,\xi_n,u_t(y,s))|
    \leq \left|\frac{\partial F^{n+1}}{\partial \xi_{n+1}} (u_{t}(x,t) -u_{t}(y,s))\right|.
    \]
    Then from Corollary \ref{c:rescale} we have 
    \[
    \begin{aligned}
    \left|\frac{\partial F^{n+1}}{\partial \xi_{n+1}} (u_{t}(x,t) -u_{t}(y,s))\right|
    &\leq  \frac{|\xi|^2 - \xi_{n+1}^2}{|\xi|^{3/2}} |u_{t}(x,t) -u_{t}(y,s))|\\
    &\leq C (|x-y|^{\beta} + |t - s|^{\beta/2})^2 (|x-y|^{\alpha} + |t - s|^{\alpha/2})r^{-\alpha- 2\gamma} \\
    &\leq C (|x-y|^{\beta} + |t - s|^{\beta/2}). 
    \end{aligned}
    \]
    We now consider the final situation in which $k=n+1$ and $j<n+1$. By multiplying the original function $u$ by $c^{-1}$, we may assume $\partial_t u \geq 1$. This multiplication will only change the estimates by a factor of $c^{-1}$. Then  since $\partial_t u \geq 1$, then $F^j$ is a convex function of $\partial_t u$, and so 
    \[
    |F^{j}(\xi_1,\ldots,\xi_n,u_t(x,t))- F^{j}(\xi_1,\ldots,\xi_n,u_t(y,s))|
    \leq \left|\frac{\partial F^{j}}{\partial \xi_{n+1}} (u_{t}(x,t) -u_{t}(y,s))\right|.
    \]
     Then 
    \[
    \begin{aligned}
     \left|\frac{\partial F^{j}}{\partial \xi_{n+1}} (u_{t}(x,t) -u_{t}(y,s)) \right|
     &\leq Cr^{\beta} r^{-\gamma} r^{-\alpha-2\gamma}(|x-y|^{\alpha}+ |t - s|^{\alpha/2}) \\
    \end{aligned}
    \]
    We just have to let $1>\beta=\alpha+3\gamma$, and the result is proven.  
    
     
\end{proof}

\end{document}
